\documentclass[11pt,a4paper]{article}
\usepackage[margin=25mm]{geometry}
\usepackage{amsmath,amssymb,amsthm,longtable,booktabs,hyperref}
\newtheorem{lemma}{Lemma}
\newcommand{\Fr}{\operatorname{Frat}}
\newcommand{\Z}{\operatorname{Z}}
\title{Checking the structural descriptions of finite hexagonal extensions}
\author{Human Compression project}\date{}
\begin{document}\maketitle
\begin{abstract}
This supplement addresses structural group assertions in the retained
source on hexagonal extensions. Exact marked presentations and
incidence computations do not by themselves identify an asserted extension.
We use their already certified permutation images to exhibit normal
subgroups, identify quotients and check the indicated splittings. No
presentation enumeration or geometric census is repeated here. Thirty-four printed shapes are verified under explicit brackets. One
split-extension assertion is disproved and replaced by a certified
nonsplit shape.
\end{abstract}
\section{What must be proved}
The source explicitly adopts the group notation of the Atlas (physical
page~3). A string containing a colon asserts a split extension; a product
sign asserts a direct product. Those assertions remain mathematical
obligations even when the action is unnamed. In the certificates below an
iterated expression is fully parenthesized. Unless parentheses in the
source dictate otherwise, extensions are read from the left; every direct
product is explicitly bracketed. We prove the resulting shape and do not
claim it uniquely determines the group among all abstract groups of that
shape.

The ambient group is not chosen by its order. We use the saved marked
rotation group of the same literal source row. For a chiral row that is
the full type-preserving group. For a regular row we use the previously
checked inversion automorphism $\alpha$ and its faithful action on two
copies: $g$ acts as $(g,\alpha(g))$ and the new involution exchanges the
copies. Its order is twice the rotation-group order. The earlier
universal-presentation and geometry certificates establish that this is
the required group; they are retained whole as premises of this
supplement.

\begin{lemma}[Structural witnesses]\label{certificate}
Let $G$ be a finite group. Suppose $N\lhd G$, with $N$ of a certified
shape $A$ and $G/N$ of a certified shape $B$. Then $G$ has extension
shape $A.B$. If moreover a subgroup $C\leq G$ satisfies
\[
 |N\cap C|=1,\qquad |N||C|=|G|,
\]
then $G=N\rtimes C$ and $C\cong G/N$. If also $[N,C]=1$, then
$G=N\times C$. These tests can be applied recursively to all the
parenthesized factors.
\end{lemma}
\begin{proof}
The quotient homomorphism has kernel $N$ and image $G/N$. Its restriction
to $C$ is injective. The order identity makes it surjective and gives
$NC=G$. Thus conjugation by $C$ supplies the action in the semidirect
product. In the commuting case multiplication $N\times C\to G$ is a
bijective homomorphism. Recursion applies these facts to the indicated
subgroups and quotient groups, rather than inferring the factors from
orders.
\end{proof}

Named simple and symmetric factors are checked by actual bijective
homomorphisms to the declared standard groups. An elementary factor
$p^a$ is checked to be elementary abelian of order $p^a$. A factor $p^{a+b}$
used here is certified with
\[
 \Z(P)=[P,P]=\Fr(P),\quad |\Z(P)|=p^a,\quad |P/\Z(P)|=p^b,
\]
with both central subgroup and quotient elementary abelian. When $a=1$
this proves that $P$ is extraspecial. The exponent and actual centre,
derived and Frattini subgroups are retained in the witness, so no stronger
choice of extraspecial type is inferred from an unsigned source symbol.
For $Q_8$ we check noncommutativity, order eight and exactly one involution;
the two nonabelian groups of order eight are then distinguished directly.
Cyclic factors are checked cyclic, not merely assigned by their order.

\section{Finite implementation and trust}
The required program reconstructs the previously accepted group and
searches normal subgroups of each prescribed order. A candidate is accepted
only after the entire kernel and quotient shapes are checked. For a colon
it then obtains an actual complement and independently checks the two
order identities in Lemma~\ref{certificate}. For a direct product it also
checks commutation on generators. The program retains the subgroup and
complement generators, not just a printed abstract description.

The normal-subgroup search can first pass to a smaller quotient. If the
required quotient has derived length at most $d$, its kernel contains
$G^{(d)}$. If its solvable radical is trivial, the kernel contains the
solvable radical of $G$. A normal $p$-subgroup lies in $O_p(G)$, so a
$p$-power kernel can instead be found among the subgroups of that core
which are normal in $G$. These are necessary conditions, not order-based
identifications; all the asserted witness properties are checked afterward.

For index two, a complement is an involution outside the kernel. It is
enough to examine a Sylow $2$-subgroup, because conjugation preserves
membership in a normal subgroup. For an $S_3$ quotient, a complement has
a subgroup of order three. Conjugate this into a Sylow $3$-subgroup and
consider its conjugacy-class representatives there. For each representative
$y$ outside the kernel, an involution $x$ inverting $y$ must lie in the
normalizer of $\langle y\rangle$. Conjugate it into a Sylow $2$-subgroup
of that normalizer. That conjugation sends $y$ to $y$ or $y^{-1}$ and
preserves the required relation. Thus these finite searches cover every
possible complement. A direct-product complement lies in $C_G(N)$; it
is a complement of $Z(N)$ there, provided $N C_G(N)=G$.

Normal-subgroup, quotient, Frattini, complement and homomorphism
computations use GAP's exact finite-group algorithms. This is a declared
computational premise, not a sample argument. A failed or timed-out search
is not counted as a structural disproof. The certificate applies only to
its exact row and literal presentation. The seven parameter-dependent or
shared L121 rows are being treated in a separate supplement, and are not
inferred from these finite rows.

\section{The correction in Table 4, row 8}
The source prints $2^{5+6}:A_5\times3:S_3:S_3$ for the row with both
parameters $(2,2)$, of full order $13\,271\,040$. Its typography was checked
directly on physical page~12. Every expression built from these factors
using only split extensions and direct products contains a subgroup
isomorphic to its $A_5$ factor, regardless of the ambiguous association.
The certified group does not contain such a subgroup.

Indeed its derived series has orders
\[
13\,271\,040,\quad3\,317\,760,\quad368\,640,\quad122\,880,
\quad1\,920,\quad120.
\]
The final subgroup $P$ satisfies $P'=P$. Every perfect subgroup $K$ of
$G$ lies in every term of the derived series: from $K\leq G^{(j)}$ we
obtain $K=K'\leq G^{(j+1)}$. Hence any $A_5$ subgroup would lie in $P$
and have index two there. An index-two subgroup is normal and gives a
nontrivial abelian quotient of $P$, contradicting $P'=P$.
As a second exact obstruction, the complete enumeration of the $120$
elements of $P$ gives exactly one involution; $A_5$ has fifteen. The
recorded order histogram is
\[
(1:1),\ (2:1),\ (3:20),\ (4:30),\ (5:24),\ (6:20),\ (10:24).
\]
This is a finite certificate for the particular exact group, not a failure
of a search for complements.

The positive corrected shape, with all association specified, is
\begin{equation}\label{corrected}
 \bigl((2^{5+6}.(A_5\times C_3)):S_3\bigr):S_3.
\end{equation}
The program exhibits a normal special subgroup of order $2^{11}$,
identifies the indicated quotient $A_5\times C_3$, and gives actual
complements at both outer colons. The inner dotted extension is nonsplit,
since a complement would again contain $A_5$. The order, exact marked
presentation and geometry of this source row are unchanged. We do not
replace the printed expression silently or claim a new historical result.

\section{Complete finite-row ledger}
The table gives the exact bracketed assertions certified here. A dot
records an extension; a colon records the actual complement checks.
Only T4.8 is a replacement of the source's split claim. The $p^{a+b}$
notation has the explicit special-group meaning stated above.
\begingroup\small
\begin{longtable}{@{}lp{133mm}@{}}
\toprule Row & Certified shape\\\midrule\endhead
T3.1 & $(((C_{2}\mathbin{.}\allowbreak \operatorname{PSL}(2,7)):\allowbreak C_{2}):\allowbreak C_{2})$\\
T3.2 & $(((C_{2}\mathbin{.}\allowbreak A_{7}):\allowbreak C_{2}):\allowbreak C_{2})$\\
T3.5 & $((Q_8\mathbin{.}\allowbreak (\operatorname{PSL}(2,7)\times\allowbreak C_{2})):\allowbreak S_{3})$\\
T4.3 & $(((3^{1+2}\times\allowbreak C_{3}):\allowbreak S_{3}):\allowbreak S_{3})$\\
T4.5 & $((3^{1+2}:\allowbreak C_{2}):\allowbreak S_{3})$\\
T4.6 & $(((2^{1+6}:\allowbreak 3^{2}):\allowbreak S_{3}):\allowbreak S_{3})$\\
T4.7 & $(3^{1+2}:\allowbreak 2^{2})$\\
T4.8 & $(((2^{5+6}\mathbin{.}\allowbreak (A_{5}\times\allowbreak C_{3})):\allowbreak S_{3}):\allowbreak S_{3})$ (corrected)\\
T4.12 & $((C_{2}\mathbin{.}\allowbreak \operatorname{PSL}(2,7)):\allowbreak S_{3})$\\
T4.14 & $((C_{3}\mathbin{.}\allowbreak \operatorname{PSU}(3,5)):\allowbreak S_{3})$\\
T4.15 & $((C_{2}\mathbin{.}\allowbreak \operatorname{PSL}(2,7)):\allowbreak C_{2})$\\
T4.16 & $((C_{2}\mathbin{.}\allowbreak \operatorname{PSL}(2,7)):\allowbreak S_{3})$\\
T4.17 & $((C_{2}\mathbin{.}\allowbreak \operatorname{PSL}(2,7)):\allowbreak C_{2})$\\
T5.3 & $(2^{1+6}:\allowbreak (A_{5}:\allowbreak C_{2}))$\\
T5.5 & $(((C_{6}:\allowbreak S_{3}):\allowbreak S_{3}):\allowbreak S_{3})$\\
T5.6 & $(((2^{1+6}:\allowbreak C_{3}):\allowbreak S_{3}):\allowbreak S_{3})$\\
T5.7 & $((((2^{1+6}:\allowbreak C_{6}):\allowbreak S_{3}):\allowbreak S_{3}):\allowbreak S_{3})$\\
T5.9 & $((((S_{3}\times\allowbreak C_{3}):\allowbreak S_{3}):\allowbreak S_{3}):\allowbreak S_{3})$\\
T5.11 & $((3^{1+2}:\allowbreak C_{6}):\allowbreak S_{3})$\\
T5.12 & $((3^{2}\times\allowbreak S_{3}):\allowbreak S_{3})$\\
T5.13 & $(((3^{1+2}\times\allowbreak C_{3}):\allowbreak S_{3}):\allowbreak S_{3})$\\
T5.14 & $((((S_{3}\times\allowbreak C_{3}):\allowbreak S_{3}):\allowbreak S_{3}):\allowbreak S_{3})$\\
T5.15 & $(((C_{4}\mathbin{.}\allowbreak (2^{4}:\allowbreak A_{5})):\allowbreak C_{2})\times\allowbreak C_{2})$\\
T5.19 & $((Q_8\mathbin{.}\allowbreak \operatorname{PSL}(2,7)):\allowbreak S_{3})$\\
T5.21 & $(((3^{2+6}:\allowbreak C_{7}):\allowbreak C_{3}):\allowbreak C_{2})$\\
T6.1 & $((2^{1+4}\times\allowbreak C_{2}):\allowbreak S_{3})$\\
T6.7 & $((2^{4}:\allowbreak S_{3}):\allowbreak S_{3})$\\
T6.9 & $(((2^{6}:\allowbreak C_{7}):\allowbreak C_{3}):\allowbreak C_{2})$\\
T7.2 & $((((2^{4}:\allowbreak A_{5}):\allowbreak C_{2}):\allowbreak C_{2}):\allowbreak S_{3})$\\
T7.5 & $((2^{9}\mathbin{.}\allowbreak \operatorname{PSL}(2,8)):\allowbreak C_{3})$\\
T7.7 & $(((2^{1+4}:\allowbreak C_{3}):\allowbreak S_{3}):\allowbreak S_{3})$\\
T8.4 & $((((2^{1+6}:\allowbreak S_{3}):\allowbreak S_{3}):\allowbreak S_{3}):\allowbreak S_{3})$\\
T8.6 & $(((2^{1+6}:\allowbreak C_{3}):\allowbreak S_{3}):\allowbreak S_{3})$\\
T8.7 & $((((2^{1+6}:\allowbreak S_{3}):\allowbreak S_{3}):\allowbreak S_{3}):\allowbreak S_{3})$\\
T8.8 & $(((C_{2}\mathbin{.}\allowbreak (A_{5}\times\allowbreak (A_{5}\times\allowbreak A_{5}))):\allowbreak C_{2}):\allowbreak S_{3})$\\
\bottomrule\end{longtable}
\endgroup

\section{Source and attribution boundary}
The source is M. E. Fernandes, D. Leemans and A. Ivi\'c Weiss, \emph{Hexagonal extensions of toroidal maps and hypermaps}, the exact retained 21-page 2018 edition. The extension descriptions themselves are source assertions. The new
work supplies their omitted explicit witnesses using the earlier saved
images; it does not claim their historical discovery. The two previously
proved order discrepancies, the local-period qualifications and all other
source dispositions remain unchanged. The bounded prior review accepts
only the exact presentations, geometry and previously verified named
groups. It is not retroactively presented as acceptance of the structures
checked here. There is no new Lean theorem, question solution or original
credit in this supplement.
\end{document}
